% Hafta 3 — AEAD (AES-256-GCM): girdiler ve çıktılar
% Derleme: py -3.12 tools/latex/derle.py docs/week-3/assets/h03-02-aead.tex
\documentclass[tikz,border=8pt]{standalone}
\usepackage{cen429-cizim}
\begin{document}
\begin{tikzpicture}[node distance=8mm and 12mm]
  % Başlık
  \node[baslik] (b) at (0,0) {AEAD (AES-256-GCM): girdiler ve çıktılar};
  \node[altbaslik] at (b.south west) {tek işlemde gizlilik + bütünlük};

  % Merkez kutu
  \node[koyukutu=30mm] (aes) at (7,-4.6) {AES-256\\GCM};

  % Sol sütun: dört girdi
  \node[kutu=32mm] (anahtar) at (0.5,-1.3) {\kb{Anahtar}\\32 bayt · gizli};
  \node[kutu=32mm] (nonce) at (0.5,-3.2) {\kb{Nonce}\\12 bayt · BENZERSİZ};
  \node[kutu=32mm] (duzmetin) at (0.5,-6.0) {\kb{Düz metin}\\korunacak veri};
  \node[kutu=32mm] (aad) at (0.5,-7.9) {\kb{İlişkili veri (AAD)}\\şifrelenmez, doğrulanır};
  \draw[ok, draw=soluk] (anahtar.east) -- (aes.north west);
  \draw[ok, draw=soluk] (nonce.east) -- (aes.west);
  \draw[ok, draw=soluk] (duzmetin.east) -- (aes.west);
  \draw[ok, draw=soluk] (aad.east) -- (aes.south west);

  % Sağ sütun: iki çıktı
  \node[iyikutu=32mm] (sifreli) at (14,-3.2) {\kb{Şifreli metin}};
  \node[iyikutu=32mm] (etiket) at (14,-6.0) {\kb{Etiket (16 bayt)}\\bütünlük kanıtı};
  \draw[iyiok] (aes.east) -- (sifreli.west);
  \draw[iyiok] (aes.east) -- (etiket.west);

  % Dip şeridi
  \node[dip, text width=160mm, below=10mm of aad.south -| aes, anchor=north] {%
    Nonce aynı anahtarla İKİ KEZ kullanılırsa gizlilik ve bütünlük birlikte çöker.};
\end{tikzpicture}
\end{document}
